\chapter{DataFrames}
\label{cap:dataframes}
% ── índice ──────────────────────────────────────────────────────
\index{DataFrame|textbf}
\index{generate@\texttt{generate}|textbf}
\index{display@\texttt{display}|textbf}
\index{filter@\texttt{filter()}|textbf}
\index{sort@\texttt{sort()}|textbf}

\section{O que é um DataFrame}

Um DataFrame é uma tabela bidimensional com colunas nomeadas. Em \hayashi{}, cada
coluna é um vetor de ponto flutuante de 64 bits. Valores ausentes são representados
por \texttt{NaN} (\textit{Not a Number}).

O DataFrame é o tipo central da linguagem — praticamente toda análise econométrica
começa com um DataFrame carregado de alguma fonte.

\section{Semântica de cópia (COW)}

DataFrames usam \textit{copy-on-write} (COW): atribuir um DataFrame a outra variável
cria um alias que compartilha os mesmos dados internos. A cópia real ocorre apenas no
momento em que um dos dois é modificado.

\begin{lstlisting}[caption={COW: baseline preservado}]
load "painel.csv" as df

let baseline = df      // alias barato, sem cópia

let df2 = df |> generate(z = x^2)    // df2 tem 'z', df não
let df3 = df |> filter(ano > 2000)   // df3 filtrado, df intacto

// baseline == df original em qualquer ponto
\end{lstlisting}

Esta semântica permite criar variantes de análise sem custo de memória até que
uma modificação real seja necessária.

\section{Criando DataFrames}

\subsection*{Bloco \texttt{input}}

Para pequenos datasets inline — ideal para exemplos e testes:

\begin{lstlisting}[caption={input block}]
input x y z
  1   2  10
  3   4  20
  5   6  30
end
\end{lstlisting}

O bloco \lstinline{input} cria um DataFrame chamado \lstinline{df} por padrão.
O número de colunas é definido pelo cabeçalho; cada linha de dados deve ter o
mesmo número de valores.

\begin{nota}
  \lstinline{input} aceita apenas valores numéricos. Use \lstinline{.} para
  representar valores ausentes. Para dados com colunas de texto, use
  \lstinline{load}.
\end{nota}

\subsection*{\texttt{dataframe(n)} — DataFrame vazio pré-dimensionado}

Cria um DataFrame vazio com um número fixo de linhas. Útil para
pré-dimensionar antes de preenchê-lo com \lstinline{generate}:

\begin{lstlisting}[caption={DataFrame pré-dimensionado}]
let df = dataframe(1000)

generate df u = rnormal(0, 1)
generate df treat = rbernoulli(0.5)
\end{lstlisting}

\subsection*{\texttt{load} — arquivos externos}

\begin{lstlisting}[caption={Carregando arquivos}]
load "dados.csv" as vendas
load "painel.dta" as painel        // formato Stata
load "resultado.parquet" as df
load "planilha.xlsx" as df         // primeira aba
load "planilha.xlsx" sheet="Dados" as df
load "base.db" as df, table=precos
load "base.db" as df, query="SELECT * FROM precos WHERE ano > 2020"
load "dados.csv" as df, sep=";"
\end{lstlisting}

Formatos suportados: CSV, TSV, JSON, Parquet, Stata (\texttt{.dta}),
Excel (\texttt{.xlsx}/\texttt{.xls}/\texttt{.ods}), SQLite, ODBC. URLs HTTP(S)
são baixadas automaticamente.

\subsection*{Projeção de colunas e filtro de linhas (\texttt{columns=}, \texttt{where=})}

Para arquivos grandes, carregar todas as colunas e linhas em RAM pode ser
desperdício ou inviável. As opções \texttt{columns=} e \texttt{where=} fazem
\textit{pushdown} da projeção e do filtro para a fonte de dados: apenas as
colunas pedidas e as linhas que casam com o predicado são materializadas em
memória.

\begin{lstlisting}[caption={Pushdown de colunas e filtro}]
load "cotacoes.parquet" as aapl, columns=[ticker, date, adj_close], where="ticker == \"AAPL\""
load "base.db" as df, table=precos, columns=[data, close], where="close > 100"
load "survey.csv" as df, where="age >= 18 && region == \"South\""
load "painel.dta" as df, columns=[id, year, y], where="year >= 2000"
load "planilha.xlsx" as df, columns=[name, score], where="score > 75"
\end{lstlisting}

\begin{center}
\begin{tabular}{llll}
\toprule
\textbf{Fonte} & \textbf{\texttt{columns=}} & \textbf{\texttt{where=}} & \textbf{Mecanismo de pushdown} \\
\midrule
Parquet   & sim & sim & Arrow \texttt{ProjectionMask} + \texttt{RowFilter} + \textit{row group pruning} por min/max \\
SQLite    & sim & sim & \texttt{SELECT cols FROM t WHERE pred} \\
ODBC      & sim & sim & idem SQLite \\
CSV / TSV & sim & sim & projeção no loop, predicado linha a linha \\
DTA       & sim & sim & projeção em \texttt{read\_record}, predicado linha a linha \\
Excel     & sim & sim & projeção após \texttt{worksheet\_range}, filtro nas células \\
JSON      & ainda não & ainda não & --- \\
\bottomrule
\end{tabular}
\end{center}

Passar \texttt{columns=} ou \texttt{where=} para um \texttt{load} JSON retorna
erro. \texttt{query=} não pode ser combinado com \texttt{columns=} ou
\texttt{where=} --- se você precisa de SQL customizado, escreva a projeção e o
filtro dentro da própria string de \texttt{query}.

\subsubsection*{Sintaxe de \texttt{columns=}}

Aceita uma lista de nomes de colunas (identificadores ou strings):

\begin{lstlisting}
columns=[ticker, date, adj_close]
columns=["ticker", "date", "adj_close"]
columns=ticker              // coluna única também aceita
\end{lstlisting}

Os nomes são validados contra o \textit{schema} da fonte antes de qualquer
dado ser lido; colunas inexistentes produzem erro listando as disponíveis.

\subsubsection*{Sintaxe de \texttt{where=}}

\texttt{where=} aceita uma expressão Hayashi da forma \texttt{coluna OP literal},
combinada com \texttt{\&\&}, \texttt{||}, \texttt{!} e \texttt{in [...]}.
A expressão é parseada pelo parser do Hayashi e normalizada em um
\texttt{RowPredicate} que cada loader avalia (ou empurra para a fonte).

\begin{lstlisting}
where="ticker == \"AAPL\""                     // igualdade de string
where="preco > 100"                            // comparação numérica
where="preco > 100 && volume > 1e6"            // AND
where="ticker in [\"AAPL\", \"MSFT\"]"         // pertinência
where="!(setor == \"Financeiro\")"             // negação
where="ano >= 2022 && produto == \"Soja\""     // combinado
\end{lstlisting}

Operadores suportados: \texttt{==}, \texttt{!=}, \texttt{>}, \texttt{>=},
\texttt{<}, \texttt{<=}, \texttt{in}, \texttt{\&\&} (e), \texttt{||} (ou),
\texttt{!} (não). Comparações devem ser \texttt{coluna OP literal} ---
comparar duas colunas não é suportado (use \texttt{generate} + \texttt{filter}
para isso).

\begin{nota}
  Em arquivos Parquet grandes, o \textit{pushdown} de filtro é especialmente
  eficaz: o predicado é avaliado durante o \textit{scan} do \textit{row group},
  antes de qualquer dado ser decodificado para a memória. Em um arquivo de
  837\,MB com 30 milhões de linhas e 8.292 \textit{tickers}, carregar
  \texttt{columns=[ticker, date, adj\_close], where="ticker == \textbackslash"AAPL\textbackslash""}
  reduziu o pico de RAM de $\sim$7{,}5\,GB (carga \textit{eager} completa) para
  $\sim$4\,MB.
\end{nota}

\subsubsection*{\textit{Row group pruning} por estatísticas (Parquet)}

Além do \texttt{RowFilter} (que avalia o predicado em cada linha durante o
\textit{scan}), o loader do Parquet lê as estatísticas de min/max de cada
\textit{row group} nos metadados do arquivo e pula \textit{row groups} inteiros
onde o predicado não pode casar. Isso é especialmente eficaz quando os dados
estão ordenados ou agrupados pela coluna filtrada.

No mesmo arquivo de 799\,MB com 30 milhões de linhas e 8.292 \textit{row groups}
ordenados por \textit{ticker}, \texttt{where="ticker == \textbackslash"AAPL\textbackslash""}
podou 8.291 de 8.292 \textit{row groups}, reduzindo o tempo de carga de
$\sim$62\,s (\textit{full scan}) para $\sim$312\,ms (212$\times$ mais rápido)
com $\sim$60\,MB de RAM.

\begin{nota}
  Para \textit{point lookups} (buscar um único valor, como um \textit{ticker}
  específico), SQLite com índice B-tree em \texttt{(ticker, date)} ainda é
  mais rápido ($\sim$42\,ms, $\sim$26\,MB de RAM) porque faz \textit{seek}
  direto sem ler metadados de todos os \textit{row groups}. Para análise
  colunar completa (todas as linhas, poucas colunas), Parquet com \textit{pruning}
  é superior devido à compressão colunar.
\end{nota}

\subsubsection*{Tipos temporais no Parquet}

Colunas do tipo \texttt{Timestamp(s|ms|$\mu$s|ns)}, \texttt{Date32} e
\texttt{Date64} no Parquet são convertidas para \textbf{strings ISO 8601}:
\texttt{YYYY-MM-DD} quando a parte de tempo é meia-noite, ou
\texttt{YYYY-MM-DDTHH:MM:SS} caso contrário. O Hayashi não possui tipo
\textit{Date} nativo, então a representação textual é a mais prática para
econometria --- permite extração de componentes via \texttt{substr} e
conversão para \textit{timestamp} Unix quando necessário.

\begin{lstlisting}[caption={Extraindo componentes de datas}]
load "cotacoes.parquet" as df, columns=[ticker, date, close]

generate df ano = substr(date, 0, 4)        // "2024"  (vetorial)
generate df mes = substr(date, 5, 2)        // "03"    (vetorial)
\end{lstlisting}

\begin{nota}
  \texttt{substr()} é vetorial e funciona dentro de \texttt{generate}, produzindo
  colunas de string. Já \texttt{date()} e \texttt{datetime()} são funções
  escalares --- aceitam uma única string ISO e retornam um \texttt{Float}
  (segundos desde 1970-01-01). Exemplo: \texttt{let ts = date("2024-03-15")}.
\end{nota}

\section{Salvando DataFrames}

\begin{lstlisting}
save df "saida.csv"
save df "saida.parquet"
save df "saida.dta"
save df "saida.xlsx"
\end{lstlisting}

\section{Inspecionando DataFrames}

\subsection*{\texttt{display}}

Exibe as primeiras linhas do DataFrame:

\begin{lstlisting}
display df
display df, 20    // exibe até 20 linhas
\end{lstlisting}

\subsection*{\texttt{describe}}

Exibe a estrutura: número de observações, variáveis e tipos:

\begin{lstlisting}
describe df
\end{lstlisting}

Saída típica:

\begin{lstlisting}[style=inline]
DataFrame: 1250 obs, 8 vars
  ano     Float  (0 missing)
  uf      Float  (0 missing)
  renda   Float  (12 missing)
  gini    Float  (3 missing)
\end{lstlisting}

\subsection*{\texttt{count} e \texttt{nrow}}

Retorna o número de linhas:

\begin{lstlisting}
let n = count(df)           // total de observações
let n2 = nrow(df)           // alias de count(df)
// contagem condicional
let n3 = count(df, renda > 1000)
\end{lstlisting}

\section{Variáveis especiais}

Dentro de comandos que operam linha a linha (\lstinline{generate}, \lstinline{mutate},
\lstinline{filter}), duas variáveis implícitas estão disponíveis:

\begin{description}
  \item[\lstinline{_n}] Número da linha atual (começa em 1).
  \item[\lstinline{_N}] Total de linhas do DataFrame.
\end{description}

\begin{lstlisting}
generate df id = _n          // 1, 2, 3, ..., N
generate df peso = 1.0/_N    // peso uniforme
\end{lstlisting}

\section{Operadores de séries temporais}

Dentro de expressões de \lstinline{generate} (forma imperativa), operadores especiais
para séries temporais estão disponíveis:

\begin{center}
\begin{tabular}{cll}
\toprule
\textbf{Operador} & \textbf{Significado} & \textbf{Exemplo} \\
\midrule
\texttt{L.x}   & Lag (valor anterior)     & \texttt{L.pib}  → pib$_{t-1}$ \\
\texttt{F.x}   & Lead (valor seguinte)    & \texttt{F.taxa} → taxa$_{t+1}$ \\
\texttt{D.x}   & Diferença primeira       & \texttt{D.pib}  → pib$_t$ - pib$_{t-1}$ \\
\texttt{L2.x}  & Lag de ordem 2           & \texttt{L2.pib} → pib$_{t-2}$ \\
\bottomrule
\end{tabular}
\end{center}

\begin{lstlisting}[caption={Operadores de séries temporais}]
generate df dpib = D.pib      // variação do PIB
generate df pib_lag = L.pib   // PIB defasado um período
\end{lstlisting}
